\documentclass[10pt,a4paper]{amsart}
\usepackage[margin=19mm]{geometry}
\usepackage{amsmath,amssymb,mathtools}
\usepackage[scr=rsfs,cal=euler]{mathalfa}
\usepackage{xfrac}
\usepackage[unicode,hidelinks,bookmarks=false]{hyperref}
\newcommand{\ie}{i.e.\ }
\newcommand{\cf}{cf.\ }
\newcommand{\resp}{respectively\ }
\newcommand{\Subsection}{\subsection}
\providecommand{\nobreakdash}{\nobreak\hbox{-}}
\setlength{\emergencystretch}{2em}
\multlinegap0pt
\usepackage{amssymb}
\usepackage{xcolor}
\usepackage[shortlabels]{enumitem}
\usepackage{tikz}
\usepackage{tcolorbox}
\usepackage{mathtools}
\newtheorem{lemma}{Lemma}
\newtheorem{proposition}{Proposition}
 \newcommand{\ZZ}{{\mathbb{Z}}}
 \newcommand{\cd}{\mathrm{cd}}
 \newcommand{\im}{\mathop{\rm Im}\nolimits}
 \newcommand{\secat}{\mathop{\rm secat}\nolimits}
 \newcommand{\tc}{\mathop{\rm tc}\nolimits}
\title{Parametrized LS category and group actions}
\author{Urban Ogrinec \and Petar Pavešić}
\date{Source: arXiv:2609.01462v1 (CC BY 4.0)}
\begin{document}
\maketitle

\noindent\emph{Source and licence.} Parametrized LS category and group actions. Version arXiv:2609.01462v1 (CC BY 4.0), distributed by arXiv under the Creative Commons Attribution 4.0 International licence, http://creativecommons.org/licenses/by/4.0/, as recorded in the arXiv licence metadata for this exact version and shown on the abstract page https://arxiv.org/abs/2609.01462v1. Full licence notice: \texttt{licenses/arxiv-2609.01462-CC-BY-4.0.txt}.\par\smallskip

\noindent\textit{Editorial scope.} Counted unit: the article's proposition Z (source0.tex lines 1015--1017). The article's complete original proof of it is delivered with it (source0.tex lines 1018--1046, one \texttt{proof} environment of 29 lines). No mathematics is written, repaired or re-derived: the statement and the proof are the article's own lines, in the article's own notation, and no retained line is substituted or re-flowed.  Retained uncounted as the source-local context the proof needs: lines 609--612.  Nothing was omitted from any retained span, so the package carries no omission record.  No retained line contains a citation, so the wrapper carries no bibliography and no bibliography entry.  No retained line contains a figure environment, an image-inclusion command or drawing code, so no image is delivered and the package declares no asset dependency.  Editorial formatting only, in addition: an AMS \texttt{amsart} preamble that restates the article's own declarations and macros used by the retained lines, namely the environments \texttt{lemma}, \texttt{proposition} on separate counters, together with the standard packages those lines need.  The retained environments print in this document as Lemma 1, Proposition 1 in order of appearance, whereas the article prints the same items with the numbering of its own full document.  The complete native source of the article as distributed in the arXiv e-print for this exact version is bundled unchanged as \texttt{sources/209114/source0.tex}.
\begin{lemma}\label{lem:secat of inclusion}
Let $H$ and $K$ be subgroups of a group $G$ such that for every $g\in G$ the conjugate 
subgroup $K^g=g K g^{-1}$ trivially intersects $H$. Then $\secat(\widetilde{BG}_K\to BG)\ge\cd(H)$.    
\end{lemma}

\begin{proposition}\label{prop:ex free by Z}
$\tc(d_w)=2$ for every $w\ne 1$.
\end{proposition}

\begin{proof}
We will base our proof on the lower estimate given in Lemma \ref{lem:secat of inclusion}.
Let $F_2$ be the free group with generators $x_1,x_2$ and let 
$\epsilon_1,\epsilon_2\colon F_2\to\ZZ$ be the homomorphisms that to each word 
$w\in F_2$ assigns the sum of exponents of $x_1$, and respectively $x_2$ in $w$.
Furthermore, let $G:=F_2\times\ZZ$, $K:=\im(s_{d_w})=\{(w^k,k)\in F_2\times\ZZ\}$ and 
$H:=\langle g\rangle \times \ZZ\le F_2\times \ZZ$, where 
$$g:=\left\{\begin{array}{ll}
x_1  & \epsilon_2(w)\ne 0 \\
x_2  & \epsilon_2(w)=0
\end{array}\right.$$
We claim that the groups $G,K,H$ satisfy the hypothesis of Lemma~\ref{lem:secat of inclusion}. 
Indeed, for elements $(u,l)\in G$ and $(w^k,k)\in K$ we have
$$(u,l)(w^k,k)(u,l)^{-1}=(uw^ku^{-1},k).$$
This is an element of $H$ if, and only if, $uw^ku^{-1}=g^l$ for some $l\in\ZZ$. 
We must consider two possibilities.

If $\epsilon_2(w)=0$, then $g=x_2$ and 
$0=\epsilon_2(uw^ku^{-1})=\epsilon_2(g^l)=l$ implies $uw^k u^{-1}=1$. In the free group
this is possible only if $w=1$, contrary to our assumption. 

If $\epsilon_2(w)\ne 0$, then $g=x_1$ and 
$k\,\epsilon_2(w)=\epsilon_2(uw^ku^{-1})=\epsilon_2(g^l)=0$ implies $k=0$. 
It follows that $l=\epsilon_1(g^l)=0$, 
therefore $H$ and the conjugate of $K$ intersect trivially. 

Finally, by Lemma~\ref{lem:secat of inclusion}, $\tc(d_w)\ge\cd(H)=2$. 
On the other hand $\tc(d_w)\le\cd(F_2\times \ZZ)=2$, therefore the equality holds.
\end{proof}

\end{document}
